\input{sol_pre.tex}
\input{q_2025-10-05_2.tex}
\begin{document}
\input{sol_head.tex}

\Savol{1}
$x$, $y$ ikkalasi toq bölsa $x+y$ juft; aks holda $x$ yoki $y$ juft. Demak $(x+y)xy$ har doim juft. (Masalan $x=1$, $y=2$ da A, B, C toq.)
\Javob{D}

\Savol{2}
$3{,}0(45)=3+\frac1{22}$, \ $2{,}0(45)=2+\frac1{22}$, \ $1{,}(08)=1+\frac8{99}$:
\[5+\frac1{11}-1-\frac8{99}=4+\frac{9-8}{99}=4{,}(01).\]
\Javob{D}

\Savol{3}
$n$ qator, har birida $n+4$ örindiq: $n(n+4)=60$ $\Rightarrow$ $n=6$.
\Javob{C}

\Savol{4}
Adras $a$, atlas $b$ (1 metr): \ $3a+5b=215\,000$, \ $0{,}88\cdot3a+0{,}92\cdot5b=194\,200$.\\
Birinçisini $0{,}92$ ga köpaytirib ayiramiz: $(2{,}76-2{,}64)a=197\,800-194\,200$ $\Rightarrow$ $a=30\,000$, $b=25\,000$.\\
Farq: $5000$.
\Javob{B}

\Savol{5}
$x^3+y^3=(x+y)(x^2-xy+y^2)$: \ $36^2-36\cdot73+73^2=3997=37^2-37\cdot73+73^2$.
\[\frac{109\cdot3997}{110\cdot3997}=\frac{109}{110}.\]
\Javob{C}

\Savol{6}
$19+6\sqrt2=\left(1+3\sqrt2\right)^2$; \ $19{,}25-1-3\sqrt2=18{,}25-3\sqrt2=\left(3\sqrt2-\frac12\right)^2$;\\
$12+2\left(3\sqrt2-\frac12\right)=11+6\sqrt2=\left(3+\sqrt2\right)^2$ $\Rightarrow$ javob $3+\sqrt2$.
\Javob{C}

\Savol{7}
$a=\sqrt[6]{26}$, \ $b=\sqrt[6]{25}$, \ $c=\sqrt[6]{27}$ $\Rightarrow$ $c>a>b$.
\Javob{C}

\Savol{8}
$a_1+8d=3a_1+6d$ $\Rightarrow$ $d=a_1$; \ $a_1+7d=2a_1+4d+8$ $\Rightarrow$ $3d-a_1=8$ $\Rightarrow$ $a_1=d=4$.\\
$S_{10}=\frac{2\cdot4+9\cdot4}2\cdot10=220$.
\Javob{C}

\Savol{9}
$b=4$, $a=4-d$, $c=4+d$: \ $(5-d)(5+d)=16$ $\Rightarrow$ $d=3$ (ösuvçi). Geometrik: $2;4;8$, $q=2$. \ $d+q=5$.
\Javob{C}

\Savol{10}
$5^{-3(a^2+4ab)}=5^{\frac43(3a^2-10ab)}$ $\Rightarrow$ $-3a^2-12ab=4a^2-\frac{40}3ab$ $\Rightarrow$ $7a^2=-\frac43ab$.\\
$a\ne0$: \ $\dfrac ab=-\dfrac4{21}$.
\Javob{D}

\Savol{11}
Kvadratlar yiğindisi $0$ $\Rightarrow$ har biri $0$: \ $a+b=5$, \ $b=2c-3$, \ $3c+a=10$.\\
$a=5-b=8-2c$ $\Rightarrow$ $3c+8-2c=10$ $\Rightarrow$ $c=2$, $b=1$, $a=4$. \ $64+1+8=73$.
\Javob{B}

\Savol{12}
Çap tomon $\pm1$: \ $x=1$ ($\sin1>0$) va $x=-1$ ($\sin(-1)<0$) — ikkalasi mos. $2$ ta.
\Javob{B}

\Savol{13}
$\cos^2t=\frac{1+\cos2t}2$: \ $\cos4x+\cos6x+\cos20x+\cos22x=0$.
\[2\cos5x\cos x+2\cos21x\cos x=0\ \Rightarrow\ 4\cos x\cos13x\cos8x=0.\]
$x=\frac\pi2+\pi k$, \ $x=\frac\pi{26}+\frac{\pi k}{13}$, \ $x=\frac\pi{16}+\frac{\pi k}8$. Eng kiçik musbat: $\frac\pi{26}$.
\Javob{B}

\Savol{14}
$t=2025^x>0$: \ $t^2-4051t+2025\cdot2026=0$ $\Rightarrow$ $t=2025$ yoki $t=2026$ — $2$ ta ildiz.
\Javob{B}

\Savol{15}
Soha: $x>-\frac32$, $x\ne-1$. \ $4x^2+12x+9=(2x+3)^2$, \ $6x^2+17x+12=(2x+3)(3x+4)$.\\
$u=\log_{3x+4}(2x+3)$: \ $2u+1+\frac1u=4$ $\Rightarrow$ $2u^2-3u+1=0$ $\Rightarrow$ $u=1$ yoki $u=\frac12$.\\
$u=1$: $2x+3=3x+4$, $x=-1$ — mos emas. \ $u=\frac12$: $(2x+3)^2=3x+4$ $\Rightarrow$ $4x^2+9x+5=0$ $\Rightarrow$ $x=-\frac54$ ($x=-1$ mos emas).\\
$1$ ta ildiz.
\Javob{A}

\Savol{16}
$x=\sqrt3$ ildiz: $3\sqrt3-3\sqrt3-3+3=0$. \ $x^3-(3+\sqrt3)x+3=\left(x-\sqrt3\right)\left(x^2+\sqrt3x-\sqrt3\right)$.\\
Kvadrat uçhad: $D=3+4\sqrt3>0$, $x=\sqrt3$ uning ildizi emas. Jami $3$ ta.
\Javob{C}

\Savol{17}
$x^2+1\ge1$. \ $x=0$ da $1\ge1$ — bajariladi. $x\ne0$ da $(x-2)(x+2)\ge12$ $\Rightarrow$ $|x|\ge4$.\\
Qabul qilmaydigan butunlar: $\pm1,\pm2,\pm3$ — $6$ ta.
\Javob{C}

\Savol{18}
$x\ne1$, $x\ne\frac12$: \ $\dfrac{(x-1)(x^2+x+1)}{(2x-1)(x-1)}=\dfrac{x^2+x+1}{2x-1}>1$ $\Rightarrow$ $\dfrac{x^2-x+2}{2x-1}>0$.\\
Surat doim musbat: $x>\frac12$, $x\ne1$.
\Javob{B}

\Savol{19}
$\sqrt x+1<3$ $\Rightarrow$ $0\le x<4$: \ $0+1+2+3=6$.
\Javob{B}

\Savol{20}
$f\left(\frac12\right)=\frac52\cdot3=\frac{15}2$, \ $f(2)=4\cdot\frac92=18$: \ $\dfrac{15}{36}=\dfrac5{12}$.
\Javob{A}

\Savol{21}
$t=\tg\frac x2$: $\cos x=\frac{1-t^2}{1+t^2}$, $\sin x=\frac{2t}{1+t^2}$, $dx=\frac{2dt}{1+t^2}$. Maxraj: $\frac{t^2-10t+25}{1+t^2}$.
\[\int_0^1\frac{2\,dt}{(t-5)^2}=\left.\frac{-2}{t-5}\right|_0^1=\frac12-\frac25=\frac1{10}.\]
\Javob{C}

\Savol{22}
$f'(x)=3x^2-2x-7$, \ $f'(2)=1=\tg\alpha$ $\Rightarrow$ $\alpha=45^\circ$.
\Javob{B}

\Savol{23}
$AO$ — bissektrisa ($O$ — içki aylana markazi). Bissektrisa xossasi: $\dfrac{AB}{AC}=\dfrac{BD}{DC}$ $\Rightarrow$ $AB\cdot DC=9\cdot4$ $\Rightarrow$ $AB^2=36$, $AB=DC=6$.\\
Bissektrisa uzunligi: $AD^2=AB\cdot AC-BD\cdot DC=54-24=30$, \ $AD=\sqrt{30}$.
\Javob{B}

\Savol{24}
$t=\cos x\in\left[-\frac12;\frac12\right]$: \ $g(t)=\dfrac{(t+3)^2}{1-t^2}$, \ $g'(t)=\dfrac{2(t+3)(3t+1)}{(1-t^2)^2}=0$ $\Rightarrow$ $t=-\frac13$.\\
$g\left(-\frac13\right)=\frac{64/9}{8/9}=8$; \ $g\left(-\frac12\right)=\frac{25}3$, $g\left(\frac12\right)=\frac{49}3$. Eng kiçigi $8$.
\Javob{C}

\Savol{25}
$OP_2=OP_3$ va burçak $60^\circ$ $\Rightarrow$ $\triangle OP_2P_3$ muntazam, $R=2$. $P_1P_2=2=R$ $\Rightarrow$ $\angle P_1OP_2=60^\circ$.\\
$P_3P_4=2\sqrt3=R\sqrt3$ $\Rightarrow$ $\angle P_3OP_4=120^\circ$, qolgani $\angle P_4OP_1=360^\circ-60^\circ-60^\circ-120^\circ=120^\circ$.
\[S=S_{OP_3P_4}+S_{OP_4P_1}=2\cdot\frac12\cdot2^2\sin120^\circ=2\sqrt3.\]
\Javob{B}

\Savol{26}
$3^2+4^2=5^2$. \ Bissektrisa: $\frac{AD}{DB}=\frac{AC}{BC}=\frac34$ $\Rightarrow$ $AD=\frac{15}7$. \ $AM=\frac52$.
\[MD=\frac52-\frac{15}7=\frac5{14}.\]
\Javob{B}

\Savol{27}
$BC=b$, $AD=a$. Örta çiziqda $KM=\frac b2$, $KN=\frac a2$ $\Rightarrow$ $MN=\frac{a-b}2=4$, $a-b=8$.\\
Teng yonli trapetsiyada diagonallar perpendikulyar bölsa, $\triangle AOD$, $\triangle BOC$ — teng yonli töğri burçakli, balandlik $h=\frac{a+b}2$.\\
$CD^2=h^2+\left(\frac{a-b}2\right)^2$ $\Rightarrow$ $240=h^2+16$, $h^2=224$, $(a+b)^2=896$.\\
$AO=\frac a{\sqrt2}$, $BO=\frac b{\sqrt2}$: \ $S_{ABO}=\frac12AO\cdot BO=\frac{ab}4$, \ $ab=\frac{(a+b)^2-(a-b)^2}4=\frac{896-64}4=208$.\\
$OK$ — $\triangle ABO$ medianasi: \ $S_{AKO}=\frac12S_{ABO}=\frac{208}8=26$.
\Javob{A}

\Savol{28}
Muntazam beşburçakda (tomoni $s$, diagonali $d$): $AG=AB=s$ ($\triangle ABG$ teng yonli), $AF=GC=d-s$.\\
$FG=d-2(d-s)=2s-d$, \ $d=\varphi s$, $\varphi=\frac{1+\sqrt5}2$: \ $FG=s(2-\varphi)=\frac s{\varphi^2}$.
\[s=2\varphi^2=3+\sqrt5,\qquad P=5\left(3+\sqrt5\right).\]
\Javob{C}

\Savol{29}
Oltiburçak balandligi (qarama-qarşi tomonlar orasidagi masofa) $a\sqrt3=6$ $\Rightarrow$ $a=2\sqrt3$.\\
$B(0;3)$, $E$ esa $B$ dan katta diagonal $2a=4\sqrt3$ masofada: $E\left(4\sqrt3;\,3\right)$. Köpaytma $12\sqrt3$.
\Javob{D}

\Savol{30}
$C$ dan $AB$ ga balandlik $CH=\frac{15\cdot20}{25}=12$. Uç perpendikulyar teoremasi: $DH\perp AB$.
\[DH=\sqrt{35^2+12^2}=37.\]
\Javob{B}

\Savol{31}
$(A\cup B)'$ — $A$ va $B$ dan taşqari: $s,f,k,g,o,x,y,j$ ($8$ ta). $C$ dan qöşimça: $m,n,l$. Jami $11$.
\Javob{D}

\Savol{32}
$x=y=1$: \ $3^5=243$.
\Javob{B}

\newpage
\Savol{33}
$\frac43\pi R^3=32\sqrt3\pi$ $\Rightarrow$ $R^3=24\sqrt3$, $R=2\sqrt3$. Silindr: radius $r$, balandlik $h$: $r^2+\frac{h^2}4=12$.

\textbf{33.} $V=\pi\left(12-\frac{h^2}4\right)h$, \ $V'=\pi\left(12-\frac{3h^2}4\right)=0$ $\Rightarrow$ $h=4$, $r^2=8$: \ $V=32\pi$ — \textbf{F}.

\textbf{34.} $S=2\pi rh=4\pi\cdot r\cdot\frac h2\le4\pi\cdot\frac{r^2+h^2/4}2=24\pi$ ($r=\frac h2=\sqrt6$ da) — \textbf{C}.

\textbf{35.} $h=4$, $r=2\sqrt2$: \ $S=2\pi\cdot2\sqrt2\cdot4=16\sqrt2\pi$ — \textbf{A}.
\Javob{33 — F, \ 34 — C, \ 35 — A}

\newpage
\Savol{36}
Ikkinçi tenglama: $\frac1y=2|x|+2$ $\Rightarrow$ $y=\dfrac1{2|x|+2}$.\\
Sistema $x\to-x$ almaştirişda özgarmaydi: $(x;y)$ yeçim bölsa, $(-x;y)$ ham yeçim. Yagona yeçim uçun $x=0$, $y=\frac12$ kerak.\\
Birinçi tenglama $x=0$ da: $\log_3(-a)=1+\log_31$ $\Rightarrow$ $a=-3$.\\
Tekşiriş ($a=-3$, $u=x^2y^2\ge0$): $x^2u+3=3(1-3u)$ $\Rightarrow$ $u\left(x^2+9\right)=0$ $\Rightarrow$ $u=0$, $x=0$. Yeçim yagona: $\left(0;\frac12\right)$.

a) $a=-3$, yiğindi $-3$.

b) $x+y=\frac12$.
\Javob{a) $-3$; \quad b) $\frac12$}

\newpage
\Savol{37}
Çap tomonni soddalaştiramiz: $\cos3x\cos x=\frac{\cos4x+\cos2x}2$, \ $\cos^2x\cos2x=\frac{\cos2x+\cos^22x}2$, \ $2\cos^4x=\frac{1+2\cos2x+\cos^22x}2$:
\[\text{Çap}=\frac{4\cos2x+2\cos^22x+1+\cos4x}2+\cos4x=1+2\cos2x+2\cos4x.\]
$\sin x\ne0$ da $\left(1+2\cos2x+2\cos4x\right)\sin x=\sin x+\left(\sin3x-\sin x\right)+\left(\sin5x-\sin3x\right)=\sin5x$.\\
($\sin x=0$, $\sin\frac x2\ne0$ da çap $=5$, öng $=\pm\frac12$ — yeçim emas.)
\[\frac{\sin5x}{\sin x}=\frac1{2\sin\frac x2}\ \Rightarrow\ \sin5x=\frac{\sin x}{2\sin\frac x2}=\cos\frac x2.\]
$5x=\frac\pi2-\frac x2+2\pi k$ $\Rightarrow$ $x=\frac\pi{11}+\frac{4\pi k}{11}$; \quad $5x=\frac\pi2+\frac x2+2\pi k$ $\Rightarrow$ $x=\frac\pi9+\frac{4\pi k}9$.

a) Eng kiçik musbat: $\frac\pi{11}$.

b) $\left(-\frac\pi2;\frac\pi2\right)$: $-\frac{3\pi}{11},\ \frac\pi{11},\ \frac{5\pi}{11}$ va $-\frac\pi3,\ \frac\pi9$ — $5$ ta.
\Javob{a) $\frac\pi{11}$; \quad b) 5}

\newpage
\Savol{38}
$t=x-1$: \ $g(t)=(t+1)^2+1$, ya'ni $g(x)=x^2+2x+2=(x+1)^2+1$.

a) Eng kiçik qiymat $1$ ($x=-1$ da).

b) $g(1)=5$, \ $f(2)=5$ $\Rightarrow$ $f^{-1}(5)=2$.
\Javob{a) 1; \quad b) 2}

\newpage
\Savol{39}
$f(0)=e$, \ $f(1)=e^{a+b+1}=e$ $\Rightarrow$ $a+b=0$. \ $f'(x)=(2ax+b)e^{ax^2+bx+1}$, $f'(0)=be=e$ $\Rightarrow$ $b=1$, $a=-1$.

a) $\dfrac ab=-1$.

b) $f(x)=e^{-x^2+x+1}$, daraja körsatkiçi $x=\frac12\in[-1{,}5;2]$ da eng katta: $-\frac14+\frac12+1=\frac54$. \ $\ln m=\frac54$.
\Javob{a) $-1$; \quad b) $\frac54$}

\newpage
\Savol{40}
Parabola $O$ dan ötadi, öqi $x=4$ $\Rightarrow$ $(8;0)$ dan ham ötadi: $f(x)=ax(x-8)$, uçi $(4;-16a)$.\\
$g$ grafigi $O$ dan ötadi, çap tarmoği koeffitsiyenti $1$: $y=x$. Uç şu töğri çiziqda: $-16a=4$, $a=-\frac14$: $f(x)=-\frac{x^2}4+2x$.\\
Öng tarmoq $y=-x+2p$ ($p=-m=n$ — siniq çiziq uçi), $(6;3)$ dan ötadi: $p=4{,}5$.

a) $n=4{,}5$, $m=-4{,}5$: \ $n-m=9$.

b) Soha: $O$ dan uçgaça $y=x$, parabola yoyi $(4;4)$ dan $(6;3)$ gaça, $(6;3)$ dan $O$ gaça $y=\frac x2$:
\[S=\int_0^4\left(x-\frac x2\right)dx+\int_4^6\left(-\frac{x^2}4+2x-\frac x2\right)dx=4+\left[-\frac{x^3}{12}+\frac{3x^2}4\right]_4^6=4+9-\frac{20}3=\frac{19}3.\]
\Javob{a) 9; \quad b) $\frac{19}3$}

\newpage
\Savol{41}
\fig{\begin{tikzpicture}[scale=0.2,font=\small]
\coordinate (B) at (0,0); \coordinate (C) at (35,0); \coordinate (A) at (12.6,16.8); \coordinate (E) at (3.75,5); \coordinate (D) at (28.33,5);
\fill[okgreen!12] (A)--(E)--(D)--cycle;
\draw (A)--(B)--(C)--cycle; \draw (E)--(D); \draw (10,5) circle (5); \draw (20,5) circle (5);
\draw[ax] (B)--(10,5); \draw[ax] (C)--(20,5); \draw[axs] (10,5)--(10,0); \draw[axs] (20,5)--(20,0);
\node[above] at (A) {$A$}; \node[below left] at (B) {$B$}; \node[below right] at (C) {$C$}; \node[left] at (E) {$E$}; \node[right] at (D) {$D$};
\node[below] at (5,0) {$10$}; \node[below] at (15,0) {$10$}; \node[below] at (27.5,0) {$15$};
\end{tikzpicture}}
Aylanalar markazlari $BC$ dan $5$ masofada, ular orasidagi masofa $10$. $B$ va $C$ dan uriniş nuqtalarigaça: $5\ctg\frac B2$ va $5\ctg\frac C2$:
\[5\ctg\frac B2+10+5\ctg\frac C2=35\ \Rightarrow\ \ctg\frac B2+\ctg\frac C2=5,\qquad \frac B2+\frac C2=45^\circ.\]
$t=\tg\frac B2$: $\tg\frac C2=\frac{1-t}{1+t}$: \ $\frac1t+\frac{1+t}{1-t}=5$ $\Rightarrow$ $6t^2-5t+1=0$ $\Rightarrow$ $t=\frac12$ (yoki $\frac13$ — simmetrik).\\
$\tg B=\frac{2\cdot\frac12}{1-\frac14}=\frac43$ $\Rightarrow$ katetlar nisbati $3:4$, $BC=35$: $AB=21$, $AC=28$.

b) $S_{ABC}=\frac12\cdot21\cdot28=294$.

a) $A$ dan $BC$ gaça balandlik $h=\frac{21\cdot28}{35}=16{,}8$. $ED\parallel BC$, $A$ dan $ED$ gaça $16{,}8-5=11{,}8$:
\[S_{ADE}=294\cdot\left(\frac{11{,}8}{16{,}8}\right)^2=294\cdot\frac{3481}{7056}=\frac{3481}{24}.\]
\Javob{a) $\frac{3481}{24}$; \quad b) 294}

\newpage
\Savol{42}
$\vec{AB}=\left(3;\,3\sqrt3\right)$, $|AB|=6$; \ $\vec{AD}=\left(1;\,-\sqrt3\right)$, $|AD|=2$.

a) $\cos\beta=\dfrac{3-9}{6\cdot2}=-\dfrac12$ $\Rightarrow$ $\beta=120^\circ$, \ $\sin\beta=\dfrac{\sqrt3}2$.

b) $BD$ diagonal böyiça bölamiz; $\vec{CB}=\left(-4;\,2\sqrt3\right)$, $\vec{CD}=\left(-6;\,-2\sqrt3\right)$:
\[S_{ABD}=\frac12\cdot6\cdot2\cdot\frac{\sqrt3}2=3\sqrt3,\qquad S_{BCD}=\frac12\left|(-4)\left(-2\sqrt3\right)-2\sqrt3\cdot(-6)\right|=\frac12\cdot20\sqrt3=10\sqrt3.\]
$S_{ABCD}=13\sqrt3$.
\Javob{a) $\frac{\sqrt3}2$; \quad b) $13\sqrt3$}

\newpage
\Savol{43}
Vektorlar bilan: $M=\frac{A+B}2$, $N=\frac{C+D}2$ $\Rightarrow$ $P=\frac{A+B+C+D}4$; \ $K=\frac{B+C}2$, $L=\frac{D+E}2$ $\Rightarrow$ $T=\frac{B+C+D+E}4$.
\[\vec{PT}=T-P=\frac{E-A}4=\frac14\vec{AE}.\]
Demak $PT\parallel AE$ va $PT=\frac{AE}4$.

a) $PT=\frac{12}4=3$.

b) $APTE$ — trapetsiya ($PT\parallel AE$), balandligi $AP\sin60^\circ=6\sqrt3$:
\[S=\frac{12+3}2\cdot6\sqrt3=45\sqrt3.\]
\Javob{a) 3; \quad b) $45\sqrt3$}

\newpage
\Savol{44}
Konus: $r=3$, $l=5$ $\Rightarrow$ $h=4$. Yarim şar radiusi $3$. Konus uçi yarim şar qutbida: silindr uzunligi $H=4+3=7$.

a) Qolgan jism sirti: silindr yon sirti, konus yon sirti va yarim şar sirti:
\[2\pi\cdot3\cdot7+\pi\cdot3\cdot5+2\pi\cdot3^2=42\pi+15\pi+18\pi=75\pi.\]

b) $V=\pi\cdot3^2\cdot7=63\pi$.
\Javob{a) $75\pi$; \quad b) $63\pi$}

\newpage
\Savol{45}
$AC=x$: \ $t(x)=\dfrac{200-x}{150}+\dfrac{\sqrt{150^2+x^2}}{75}$.
\[t'(x)=-\frac1{150}+\frac x{75\sqrt{150^2+x^2}}=0\ \Rightarrow\ 2x=\sqrt{150^2+x^2}\ \Rightarrow\ x=50\sqrt3.\]
$CT=100\sqrt3$, $BC=200-50\sqrt3$.

a) $t=\dfrac{200-50\sqrt3}{150}+\dfrac{100\sqrt3}{75}=\dfrac43-\dfrac{\sqrt3}3+\dfrac{4\sqrt3}3=\dfrac43+\sqrt3\approx3{,}07\approx3$ daqiqa.

b) $BC+CT=200-50\sqrt3+100\sqrt3=200+50\sqrt3\approx200+85=285$ m.
\Javob{a) 3 daqiqa; \quad b) 285 m}

\ShowJavoblar
\end{document}
